\subsection{The involution question and its precise realization obstruction}

\begin{result}{B.3}{Open}{Conjugacy of all GT involutions to complex conjugation}
The assertion that every involution in $\GT$ is conjugate to
$c=(-1,1)$ in $\GT$ remains unresolved in this audit.
The supplied, refereed torsion theorem says that every
nonidentity finite-order element has order two and multiplier
$-1$. The checked theorem of Lochak--Schneps gives
$C_{\GT}(c)=\langle c\rangle$. Neither assertion supplies
a conjugator for an arbitrary involution.
\end{result}

Here is the exact comparison for the attempted E.44 method.
That method starts with an actual finite subgroup
$H\hookrightarrow\widehat\Gamma_{0,n}$. Boggi--Zalesskii
geometrizes that subgroup and identifies its normalizer and
centralizer through the completion map from the discrete
mapping class group. An involution $t\in\GT$ instead
gives an outer automorphism
$t\in\operatorname{Out}(\widehat F_2)$.
The canonical map
\[
 \widehat F_2\xrightarrow{\operatorname{Inn}}
 \operatorname{Aut}(\widehat F_2)
 \longrightarrow\operatorname{Out}(\widehat F_2)
\]
is the zero homomorphism. On the other hand the action of
$t$ on $\widehat F_2^{\rm ab}$ is $-I$, nontrivial already
modulo three. Thus this particular map cannot realize $t$
as the subgroup required by that argument. This proves
the failure of the proposed direct substitution; it is
not a nonconjugate GT involution.

One can form the genuine semidirect extension
$\widehat F_2\rtimes\langle t\rangle$ after choosing the
normalized representative of the involution. To obtain
a conjugator inside $\GT$, the missing assertion would
be a realization and comparison theorem for that specific
extension preserving all three peripheral classes and
the five-point structure. A normalizer in another
completed surface group still needs a map with these
properties. The supplied LM4\_SUCCESSOR A.39 reduces
the required conjugator to the specified unit $\mu$
satisfying its three-peripheral equation; its A.81 gives
an explicit counterexample to the proposed automatic
lifting principle. We retain that reduction without
repeating it or using the disproved lifting principle.
The smallest undecided case here is one arbitrary
multiplier $-1$ GT involution with a nontrivial normalized
parameter $f$, for which no such $\mu$ has been produced.

The primary assertion used above is precisely
\cite[Proposition 4(ii), proof on printed p.~581]{CA_LS1997}.
The proof computes the centralizer of the distinguished
involution by its cohomological comparison; its conclusion
is the displayed centralizer equality. We use its
inclusion of the two known elements $1,c$ and its
exclusion of every other commuting element. The
existence of a conjugator from a second involution
would require an additional assertion about another
orbit, and is the unresolved step recorded here.

We also checked two possible recognition inputs beyond that
centralizer theorem. Herfort--Zalesskii's pro-$p$ recognition
theorem \cite[Theorems 6.2--6.3]{HerfortZalesskii2013}
applies to the image of the extension under the actual
completion map
\[
 \widehat F_2\rtimes\langle t\rangle
       \longrightarrow F_2^{(p)}\rtimes\langle t_p\rangle.
\]
The supplied LM4\_SUCCESSOR A.42 already constructs the
prescribed conjugators at every coefficient prime, so this
recognition input does not resolve the remaining full
third-peripheral equation of A.39. We retain that established
local progress and have not repeated it as a new result.

The 2026 theorem of Wykowski
\cite[Theorem 1.2]{Wykowski2026} starts with a finitely
generated discrete virtually free group $\Gamma$ and a
specified finitely generated subgroup $\Lambda\leq\Gamma$.
It compares a factorization of $\widehat\Gamma$ containing
$\overline\Lambda$ with a discrete amalgamated factorization.
For the extension attached to an arbitrary $t$, this attempted
application still lacks a discrete $\Gamma$ and a comparison
isomorphism
$\widehat\Gamma\simeq\widehat F_2\rtimes\langle t\rangle$
retaining the three peripheral groups. Its Theorem 1.1
constructs a derivation from an already specified accessible
factor; it does not construct the needed factor from the
augmentation module alone. These are the exact unverified
inputs of this recognition route.

\subsection{Prime-power quotients and a possible arithmetic separator}

\begin{result}{B.4}{Open}{The universal quotient assertion and a separator from the Galois image}
The quotient comparisons in A.1 and A.3 prove the vanishing
of their specified CM discrepancy on $\GT$.
B.1.7 and B.4.1 further identify the original GT parameter
on every genus-zero prime-power quotient with all branch
points marked and the stated regular-orbit markings.
The universal assertion over all genera and all frame
characters remains unresolved, as does a continuous
invariant separating $G_{\mathbb Q}$ from $\GT$.
\end{result}

The actual positive construction is the restriction of
the tower automorphism to the completed centralizer of
the chosen finite subgroup, followed by the homomorphism
to the mapping class group of its quotient orbifold.
To identify that quotient's parameter with the original
parameter, A.1 uses two named forgetful maps and
an isomorphism of their marked source; A.3 uses four
such maps and the classification of the eight-point
outer automorphism group. For cyclic sphere actions,
B.1.7 supplies the required choice subgroups and
power-cover/opening square. B.4.1 gives the corresponding
dihedral normalizer factorization, completing this
genus-zero comparison. For a higher-genus prime-power
action, the finite-subgroup theorem gives its geometric
realization and completed centralizer; parameter-detecting
comparison maps with the required branch and handle
markings remain to be constructed. The prime-power
theorem alone contains no assertion that the resulting
parameter is the original one.

The unproved universal phrase \emph{no frame character}
also quantifies over characters beyond the two
specified signed frames. Our exact calculation shows
that the particular homomorphism $\Delta$ restricts to
zero on $\GT$, and R1 shows its restriction to the
arithmetic image is zero. This is the restriction
map on one named character. The question about every
other character requires its definition and its own
restriction map. We have not inferred it from the
vanishing of this one restriction.

The attempts examined here are the quotient centralizers,
the harmonic covering comparison, the literal descent
operator of B.1, and the section constructions of
Part C. The arithmetic descent formula in B.1
identifies a normalization effect and therefore cannot
by itself separate the two groups. The fixed-support
and fixed-ramification section theorems in Part C
do not cover the infinite-support ordinary pairs
of the prompt. A separator from any of these routes
has not been obtained. The inclusion
$G_{\mathbb Q}\hookrightarrow\GT$ used throughout is
the checked Drinfeld--Ihara map on the tangential
free fundamental group; deciding its surjectivity
remains beyond the results established here.
